\documentclass[10pt,a4paper]{amsart}
\usepackage[margin=19mm]{geometry}
\usepackage{amsmath,amssymb,mathtools}
\usepackage[scr=rsfs,cal=euler]{mathalfa}
\usepackage{xfrac}
\usepackage[unicode,hidelinks,bookmarks=false]{hyperref}
\newcommand{\ie}{i.e.\ }
\newcommand{\cf}{cf.\ }
\newcommand{\resp}{respectively\ }
\newcommand{\Subsection}{\subsection}
\providecommand{\nobreakdash}{\nobreak\hbox{-}}
\setlength{\emergencystretch}{2em}
\multlinegap0pt
\usepackage[shortlabels]{enumitem}
\usepackage{bm}
\usepackage{amssymb}
\usepackage[normalem]{ulem}
\usepackage{xcolor}
\newtheorem{assumption}{Assumption}
\newtheorem{proposition}{Proposition}
\newcommand{\R}{\mathbb{R}}
\newcommand{\bG}{\mathbf{G}} %grafico
\DeclareMathOperator{\dist}{dist}
\newcommand{\norm}[1]{\left\lVert#1\right\rVert}
\title{Stationary and stable varifolds with singularities}
\author{Camillo De Lellis, Jonas Hirsch, Luca Spolaor}
\date{Source: arXiv:2509.21508v2 (CC BY 4.0)}
\begin{document}
\maketitle

\noindent\emph{Source and licence.} Stationary and stable varifolds with singularities. Version arXiv:2509.21508v2 (CC BY 4.0), distributed by arXiv under the Creative Commons Attribution 4.0 International licence, http://creativecommons.org/licenses/by/4.0/, as recorded in the arXiv licence metadata for this exact version and shown on the abstract page https://arxiv.org/abs/2509.21508v2. Full licence notice: \texttt{licenses/arxiv-2509.21508-CC-BY-4.0.txt}.\par\smallskip

\noindent\textit{Editorial scope.} Counted unit: the article's proposition prop:localglue (source0.tex lines 298--305). The article's complete original proof of it is delivered with it (source0.tex lines 384--391, one \texttt{proof} environment of 8 lines, which the article places away from the statement). No mathematics is written, repaired or re-derived: the statement and the proof are the article's own lines, in the article's own notation, and no retained line is substituted or re-flowed.  Retained uncounted as the source-local context the proof needs: lines 261--267; lines 276--278.  Nothing was omitted from any retained span, so the package carries no omission record.  No retained line contains a citation, so the wrapper carries no bibliography and no bibliography entry.  No retained line contains a figure environment, an image-inclusion command or drawing code, so no image is delivered and the package declares no asset dependency.  Editorial formatting only, in addition: an AMS \texttt{amsart} preamble that restates the article's own declarations and macros used by the retained lines, namely the environments \texttt{assumption}, \texttt{proposition} on separate counters, together with the standard packages those lines need.  The retained environments print in this document as Assumption 1, Proposition 1 in order of appearance, whereas the article prints the same items with the numbering of its own full document.  The complete native source of the article as distributed in the arXiv e-print for this exact version is bundled unchanged as \texttt{sources/185749/source0.tex}.
\begin{assumption}\label{ass:glue}
Let $u_i\colon \Omega_{7,12}\to \R$ be two minimal graphs that satisfy
\begin{enumerate}[label=(A.\arabic*)]
    \item\label{assumption.graph1} there is $x_0 \in \Omega_{7,12}$ with $u_1(x_0)=u_2(x_0)=0$
    \item\label{assumption.graph2} $\norm{\nabla u_1}_{L^\infty(\Omega_{7,12})}+\norm{\nabla u_2}_{L^\infty(\Omega_{7,12})} = \boldsymbol{\delta}<\boldsymbol{\delta}_0$.
\end{enumerate}
\end{assumption}

\begin{equation}\label{eq.gluefunct}
v(x)=(1-\theta_9(|x|))\, u_1(x)+\theta_9(|x|)\, u_2(x)\,.
\end{equation}

\begin{proposition}[Local gluing]\label{prop:localglue}
    Let $u_1,u_2$ be as in Assumptions \ref{ass:glue} and let $v$ be as in \eqref{eq.gluefunct}. Then there exists a \emph{closed} form $\omega$ defined in a neighborhood of the graph of $v$ such that
    \begin{enumerate}[label=(F.\arabic*)]
    \item\label{prop:localglue-F1} $\omega|_{\bG_v}=d\operatorname{vol}_v$, that is $\omega$ is identically one on the tangent to $\bG_v$;
    \item\label{prop:localglue-F2} $\omega=\omega_{u_1}$ on $\Omega_{7,8}\times [-10,10]$ and $\omega=\omega_{u_2}$ on $\Omega_{11,12}\times [-10,10]$;
    \item\label{prop:localglue-F3} $\|\omega\|-1 = \dist(y,\bG_v)\,R$ with $\norm{D^k R}_{L^\infty(\Omega_{7,12}\times[-10,10)} \lesssim_k \,\boldsymbol{\delta}$ for any $k$.
    \end{enumerate}
\end{proposition}

\begin{proof}
    %Recall that $\norm{w}=1$ on $\bG_v$, compare \ref{prop:localglue-F1} and $\norm{w(x)}=1$ for all $x \notin \Omega_{8,11}\times \R$. Since 
    %\[\norm{D(\norm{w}-1)}_{L^\infty} \lesssim \norm{Dw}_{L^\infty}\lesssim \boldsymbol{\delta}\]
    %we deduce that the same holds true for $k=0$ i.e. $\norm{\norm{w}-1}_{L^\infty}\lesssim %\boldsymbol{\delta}$ and therefore for all $k$. 
    For a smooth $\theta$ with $\theta(t) \equiv 1$ for $t< 2$ and $\theta(t)=0$ for $t>3$, the interpolation
    \[ \lambda(x)=1+\theta(\dist(y,  {\rm spt}\, (\bG_v)))\left(\norm{\omega(y)}-1\right)\]
    has the desired properties, by Proposition \ref{prop:localglue}.
\end{proof}

\end{document}
